\documentclass[11pt]{article}
\usepackage[a4paper, margin=2.5cm]{geometry}
\usepackage{amsthm}
\usepackage{graphicx}
\usepackage[slovene]{babel}
\usepackage[T1]{fontenc}
\begin{document}
\author{Matej Rojec}
\title{Brownovo gibanje}
% Brownovo gibanje
% Matej Rojec
\maketitle
\newtheorem{definicija}{Definicija}
\newtheorem{izrek}{Izrek}
\newcommand{\f}{{F}}


Brownovo gibanje (več v \cite{thm:stopped_brownian}) je intuitivno slučajen proces, % Sklic na knjigo
ki predstavlja naključno gibanje delcev v mediju.
    
    % Slika: PerrinPlot2.pdf, velikost scale=0.3
    % Napis pod sliko: 
    % Reprodukcija slike iz Jean Baptiste Perrin, \emph{Mouvement brownien et réalité moléculaire}, Ann. de Chimie et de Physique (VIII) 18, 5-114, 1909
    \begin{figure}[htp]
        \centering
        \includegraphics[scale=0.3]{PerrinPlot2.pdf}
        \caption{Reprodukcija slike iz Jean Baptiste Perrin, \emph{Mouvement brownien et réalité moléculaire}, Ann. de Chimie et de Physique (VIII) 18, 5-114, 1909}
        \label{fig:PerrinPlot2}
    \end{figure}
    % Začetek definicije
    
    \begin{definicija}
        Standardno Brownovo gibanje \(\{B_t\}_{t \geq 0}\) je slučajen proces z naslednjimi lastnostmi: 
            
            \begin{enumerate}
                \item \(B_0 = 0\).
                \item Prirastki \(B_{t_n} - B_{t_{n-1}}, B_{t_{n-1}} - B_{t_{n-2}}, \ldots, B_{t_2} - B_{t_1}, B_{t_1} - B_{t_0}\) so neodvisne slučajne spremenljivke, za vsak \(t_0 \leq t_1 \leq \cdots \leq t_{n-1} \leq t_n\).
                \item Za vsak \(t \geq 0\) in \(h > 0\) velja \(B_{t+h} - B_t \sim \mathcal{N}(0, h)\).
                \item Funkcija \(t \mapsto B_t\) je zvezna skoraj gotovo.
            \end{enumerate}
    \end{definicija}
    % Konec definicije
    
    Preden zapišemo izrek, definirajmo še pojem časa ustavljanja.
    
    % Začetek definicije
    
    \begin{definicija}
        Slučajna spremenljivka \(\tau\) na verjetnostnem prostoru \(\f\) z vrednostmi v \(\f\)
        je čas ustavljanja glede na filtracijo \(\f\), če velja: \(\f\).
    \end{definicija}
    % Konec definicije
    
    Zdaj lahko zapišemo izrek ??. % Sklic na izrek z oznako thm:stopped_brownian
    
    % Začetek izreka
    
    \begin{izrek}
        Naj bo \(\{B_t\}_{t \geq 0}\) (standardno) Brownovo gibanje, \(\tau\) čas ustavljanja glede na
        \(\f\) in naj velja \(\f\).
        Potem je tudi proces:
        \[
        \hat{B} := \{B_{\tau+t} - B_\tau \mid t \geq 0\}
        \]
        (standardno) Brownovo gibanje in neodvisen od \(\f\).
    \end{izrek}
    % Konec izreka
    \bibliographystyle{siam}
    \bibliography{magic}
    
\end{document}